\section{Introduction}

Despite decades of progress, protein folding remains challenging to mechanistically model because the search space of 3D conformations grows combinatorially with chain length and because physically meaningful pathways are hard to extract from large-scale sampling. Modern structure predictors---notably \texttt{AlphaFold2} (\texttt{AlphaFold}) and RoseTTAFold---can recover native-like folds for many sequences, but they typically provide limited interpretability for \emph{why} a trajectory moves through certain intermediate states, and they rarely provide a mechanism-level, replayable account of how constraints are enforced during folding. Molecular dynamics provides mechanistic trajectories, but at a computational cost that is prohibitive for broad exploration, and the resulting trajectories are difficult to summarize into compact, auditable decision programs.

We refer to our framework as \texttt{PhasonFold}.

PhasonFold takes a different approach: treat folding as a \textbf{discrete encoding} with auditable dynamics. A fold is not merely a final 3D structure; it is a \emph{move sequence} executed in a discrete state space, with explicit feasibility checks and \emph{auditable logs}. This philosophy motivates three design choices:

\paragraph{(i) Discrete state and move algebra.}
We represent the folding state as a compact discrete state vector (a double-rail encoding), enabling discrete moves and deterministic replay.

\paragraph{(ii) Certificate-driven nonlocal moves.}
We curate a bank of nonlocal basin-hopping moves and apply them adaptively based on a frustration/roughness proxy derived from observed residual geometry.

\paragraph{(iii) Geometry via 6D cut-and-project and phason feasibility.}
We link the discrete walk to 3D coordinates through a 6D icosahedral cut-and-project map. This exposes a perpendicular-space deviation—a \emph{phason} variable—that can be used as a soft shaping term or as a feasibility gate. Crucially, we show that phason becomes practically useful when implemented as a \emph{projector} (repair) rather than a pure filter.

PhasonFold adopts this stance: it represents folding as a replayable computation over a discrete state vector and encodes trajectories as sequences of local and nonlocal moves that can be replayed and audited. At the same time, our geometric hypothesis suggests that biologically realized backbones may concentrate near a low-phason submanifold induced by a 6D icosahedral cut-and-project geometry. This remains a testable hypothesis under our current operational proxy; accordingly, we treat phason not merely as a post hoc descriptor but as an executable dynamical feasibility variable that can guide and relax trajectories by projecting states back toward an acceptance window.

For additional conceptual context on the geometric viewpoint (information geometry, geodesics, and priors), see \texttt{docs/theory/protein\_folding\_geometry.md}.
